% sections/ndss_01_introduction.tex
% NDSS 2027 — CapMismatch Introduction
% Target: 1.5 columns / ~750 words

The emergence of LLM agent skill marketplaces---platforms such as
\clawhub (\NskillsCH skills, $>$2.1M installs) and skills.sh that let
developers publish installable tools extending language model
capabilities---represents a new class of software supply chain~\cite{xi2023rise,schick2024toolformer}.
Unlike traditional package managers~\cite{zimmermann2019small},
skills operate with \emph{NLP-driven ambient authority}:
once installed, a skill executes on behalf of the user,
reading files, executing shell commands, accessing credentials,
and sending network requests, often without per-action confirmation
or OS-level sandboxing.

Critically, the \emph{only} capability disclosure mechanism available
to users and platform operators is the skill's natural-language description.
There is no machine-readable permission manifest, no capability whitelist,
and no runtime sandbox to fall back on when the description is inaccurate.
This creates a fundamental \emph{compositional security failure}:
when users combine multiple skills---a web-search source with a
development-utility sink, an email fetcher with a memory manager---
the aggregate set of capabilities the agent can exercise
may far exceed anything disclosed in any individual skill's description.
Neither the user nor the platform can anticipate the attack surface
created by each unique skill combination.

\paragraph{Our contributions.}
We make four contributions:

\begin{enumerate}[leftmargin=*, label=\textbf{C\arabic*.}]

\item \textbf{Measurement at scale.}
We analyze all \NskillsCH skills in \clawhub using
\lscps (Layered Structured Capability Pattern Search),
a rule-based static analysis pipeline for extracting security-sensitive
capabilities from skill implementation code.
We find that \ICone of skills with dangerous capabilities
\emph{completely omit} them from their visible descriptions
(mean Capability Recall $\overline{\text{CR}} = \MeanCR$),
constituting an ecosystem-wide transparency failure.

\item \textbf{IC taxonomy.}
We introduce the Inconsistency Class (IC) taxonomy---IC-0 through IC-4---that
characterizes four structural patterns of capability-description mismatch,
from full omission (IC-1, 75.8\%) to false claims (IC-3, 0.4\%).
This taxonomy provides the first systematic vocabulary for reasoning
about skill transparency.

\item \textbf{Compositional attack surface validation (CAPE).}
We demonstrate that capability opacity does not merely create theoretical risk;
it creates a \emph{measurable, exploitable compositional attack surface}.
When any source skill (web search, email, calendar, file reader, memory,
browser automation, database connector) is co-installed with an IC-1 sink
skill, the combination constitutes an attack surface: an adversary
who can influence the content delivered by the source skill can trigger
the sink's undisclosed dangerous capability without any
additional platform access.
This attack surface arises from the composition itself---
neither source nor sink is individually malicious,
but together they expose capabilities the user never consented to.

We operationalize this surface through \cape (Compositional Attack Pair Evaluator),
testing \NpairsTotal unique source-sink pairs spanning seven injection-channel
categories (C1: LLM context, C2: filesystem, C3: agent memory, C5: database,
C9: boot hook / cron, C13: network) on a real OpenClaw /
\texttt{gpt-5.3-codex} agent platform.
Each source skill is individually registered as an
identity-spoofing MCP server mirroring the real skill's name and tool schema,
modeling the realistic supply-chain threat of a malicious clone skill.
We report two pre-defined metrics:
\textbf{ITT}: \NpairsSuccessITT/\NpairsTotal = \ITTrate\% (95\% CI: \ITTlo--\ITThi\%),
counting infrastructure timeouts as failures;
\textbf{completed-run}: \NpairsSuccessCR/\NpairsCRDenom = \CRrate\%
(95\% CI: \CRlo--\CRhi\%)
conditioned on pairs where the agent returned a response.
Side effects include shell command execution, file deletion,
cron backdoor installation, cross-session memory poisoning,
process termination, wallet exfiltration, and privilege escalation.
33/33 tested shell-execution sink pairs succeed (completed-run: 100\%),
revealing \emph{no consistent runtime defense} against this attack class.

\item \textbf{Temporal analysis.}
We track 24,248 version transitions in \clawhub and find that dangerous
capabilities are silently introduced in updates with no description change (\Tsixrecall recall at Tier~6).
Platform scanners detect 0\% of these transitions.

\end{enumerate}

\paragraph{Scope and novelty vs. prompt injection.}
Our work is related to but distinct from prompt injection
attacks~\cite{greshake2023indirect,perez2022ignore}.
Prompt injection studies \emph{the manipulation of individual inputs}
to produce unsafe agent behavior; we study
\emph{the compositional attack surface created by installing
combinations of skills with undisclosed capabilities}.
The distinction is architectural:
our attack requires no crafted user input and no runtime adversarial access.
The attack surface exists the moment a dangerous IC-1 sink skill
is installed alongside any source skill---it is a \emph{structural}
vulnerability of the skill composition, not an input-validation failure.
This means existing prompt-injection defenses (input sanitization,
constrained decoding, system-prompt hardening) are \emph{not applicable}:
the dangerous capability $\kappa$ is invoked through the agent's
normal, intended tool-use workflow, with attacker-influenced content
arriving through a legitimate, already-trusted source channel.

\paragraph{Ethical considerations.}
All experiments were conducted in fully isolated environments
with no impact on real users, external services, or production systems.
Skill identifiers used in this paper are real ClawHub skill names
included to enable reproducibility; no exploit code is released.
We have privately notified the ClawHub platform operators of our findings.

\paragraph{Paper organization.}
\S\ref{sec:background} presents the threat model.
\S\ref{sec:taxonomy} introduces the IC taxonomy.
\S\ref{sec:measurement} reports the large-scale measurement.
\S\ref{sec:cape} presents the \cape attack validation.
\S\ref{sec:temporal} analyzes version-update risks.
\S\ref{sec:discussion} discusses limitations and defenses.
\S\ref{sec:related} situates our work.
